\documentclass[10pt,a4paper]{amsart}
\usepackage[margin=19mm]{geometry}
\usepackage{amsmath,amssymb,mathtools}
\usepackage[scr=rsfs,cal=euler]{mathalfa}
\usepackage{xfrac}
\usepackage[unicode,hidelinks,bookmarks=false]{hyperref}
\newcommand{\ie}{i.e.\ }
\newcommand{\cf}{cf.\ }
\newcommand{\resp}{respectively\ }
\newcommand{\Subsection}{\subsection}
\providecommand{\nobreakdash}{\nobreak\hbox{-}}
\setlength{\emergencystretch}{2em}
\multlinegap0pt
\usepackage{bm}
\usepackage{bbm}
\usepackage{dsfont}
\usepackage{xspace}
\usepackage{cancel}
\usepackage{mathtools}
\usepackage{amsthm}
\makeatletter
\@ifundefined{lemma}{\newtheorem{lemma}{Lemma}}{}
\makeatother
\title[Lie symmetry classification and invariant solutions of time-]{Lie symmetry classification and invariant solutions of time-fractional telegraph systems with variable coefficients}
\author{Sodbaatar Adiya, Khongorzul Dorjgotov, Bayarmagnai Gombodorj, Bayarpurev Mongol, Uuganbayar Zunderiya}
\date{Source: arXiv:2604.25079v1 (CC BY 4.0)}
\begin{document}
\maketitle

\noindent\emph{Source and licence.} Lie symmetry classification and invariant solutions of time-fractional telegraph systems with variable coefficients. Version arXiv:2604.25079v1 (CC BY 4.0), distributed under the Creative Commons Attribution 4.0 International licence, as recorded in the arXiv licence metadata for this exact version and shown on the abstract page https://arxiv.org/abs/2604.25079v1. Full rights notice: licenses/arxiv-2604.25079-CC-BY-4.0.txt.\par\smallskip

\noindent\textit{Editorial scope.} Counted unit: the Lemma whose source label is \texttt{lemma:linear\_inf} in the article (source0.tex lines 291--293, source label \texttt{lemma:linear\_inf}), delivered together with the article's own complete original proof of it (source0.tex lines 295--329, one \texttt{proof} environment of 35 lines). No mathematics is written, repaired or re-derived: statement and proof are the article's own text line for line. Required context retained uncounted: duruv (lines 264--275); tav (lines 264--275); zurgaa (lines 264--275); arav (lines 277--279). Every definition, display and label of those lines is retained unchanged. Omitted from this package and not redistributed: 1--263, 276--276, 280--290, 294--294, 330--823. Nothing inside a retained excerpt is omitted or altered: every retained body is delivered exactly as the article's lines are written, so no omission receipt of any kind accompanies this package. Citations: the retained excerpts carry no citation command at all, so no bibliography entry of the article is transcribed and no reference is redistributed. Reference adaptation: none was needed and none was made; every reference inside the delivered unit resolves inside it, so no unresolved reference or citation is produced. Printed slips kept literal: no printed word, symbol, hypothesis or number of the counted statement or of its proof was altered. Layout and class: the article's own class is \texttt{amsart} with packages including amssymb, latexsym, amsmath, amsthm, enumitem, mathrsfs, geometry, fullpage, fancyhdr, xcolor, bbm, stmaryrd, hyperref, lineno; the wrapper is the lane's pinned \texttt{amsart} editorial wrapper at 10pt on A4, which restates the theorem-like environments the retained text needs and adds the article's own macro definitions that the retained text needs (none), with extra wrapper packages (bm, bbm, dsfont, xspace, cancel, mathtools). The delivered numbering is the unit's own. Assets: no retained span contains a figure environment, a graphics-inclusion command, a PDF-page inclusion, a table or a TikZ picture; no image file is copied into this package, the package declares no asset dependency and no image file is redistributed with it. Licence: version arXiv:2604.25079v1 of this article is distributed under the Creative Commons Attribution 4.0 International licence, as recorded in the arXiv licence metadata for this exact version and shown on the abstract page https://arxiv.org/abs/2604.25079v1. A full rights notice is delivered at licenses/arxiv-2604.25079-CC-BY-4.0.txt. Characters: every retained excerpt is pure ASCII, so the delivered engine, XeLaTeX with its default Latin Modern text font, reports no missing character.
\begin{eqnarray}
& & \tau_x=\tau_u=\tau_v=\xi_u=\xi_v=0,\label{aneg}\\
& & D_t^n(\xi)=0,\qquad n=1,2,\ldots,\label{gurav}\\
& & \mu_u-\alpha D_t(\tau)-\phi_v+\xi_x=0,\label{duruv}\\
& & f(x)\phi_v-\alpha f(x) D_t(\tau)-f(x)'\xi-f(x)\mu_u+f(x)\xi_x=0,\label{es}\\
& & \binom{\alpha}{n}\frac{\partial^n\mu_u}{\partial t^n}-\binom{\alpha}{n+1}D_t^{n+1}(\tau)=0,\qquad n=1,2,\ldots,\label{neg}\\
& & \frac{\partial^n\mu_v}{\partial t^n}=0,\qquad n=1,2,\ldots,\label{hoyor}\\
& & f(x)\mu_v-\phi_u=0,\label{tav}\\
& & \frac{\partial^n\phi_u}{\partial t^n}=0,\qquad n=1,2,\ldots,\label{doloo}\\
& & \binom{\alpha}{n}\frac{\partial^n\phi_v}{\partial t^n}-\binom{\alpha}{n+1}D_t^{n+1}(\tau)=0,\qquad n=1,2,\ldots,\label{naim}\\
& & \frac{\partial^\alpha\mu}{\partial t^\alpha}-u\frac{\partial^\alpha\mu_u}{\partial t^\alpha}-v\frac{\partial^\alpha\mu_v}{\partial t^\alpha}+g(x)u\mu_v-\phi_x=0,\label{zurgaa}
\end{eqnarray}

\begin{equation}\label{arav}
\frac{\partial^\alpha\phi}{\partial t^\alpha}-v\frac{\partial^\alpha\phi_v}{\partial t^\alpha}-u\frac{\partial^\alpha\phi_u}{\partial t^\alpha}+g(x)u\phi_v-\alpha g(x)uD_t(\tau)-g(x)'u\xi-g(x)\mu-f(x)\mu_x=0.
\end{equation}

\begin{lemma} \label{lemma:linear_inf}
 $\mu(x,t,u,v)$ and $\phi(x,t,u,v)$ are linear functions with respect to $u$ and $v$. Thus, $\mu_1=0$ and $\phi_1=0$.
\end{lemma}

\begin{proof}
From \eqref{duruv} and \eqref{tav} by differentiating them with respect to $u$ and $v$ we obtain
\begin{align}
\mu_{uu} - \phi_{vu} &= 0, \label{11u} \\
\mu_{uv} - \phi_{vv} &= 0, \label{11v} \\
f(x)\mu_{vu} - \phi_{uu} &= 0,\label{15u}\\
f(x)\mu_{vv} - \phi_{uv} &= 0 \label{15v}.
\end{align}
By subtracting \eqref{15v} from \eqref{11u}, and by subtracting \eqref{15u} from the product of $f(x)$ and \eqref{11v},  we arrive at the following relations
\begin{align}
\mu_{uu} -  f(x) \mu_{vv} &= 0, \label{11u15v} \\
\phi_{uu} - f(x)\phi_{vv} &= 0 \label{f11v15u}. 
\end{align}
By taking second-order derivative from \eqref{zurgaa} with respect $u$ and $v$ we obtain
\begin{align}
-\frac{\partial^\alpha\mu_{uu}}{\partial t^\alpha}-u\frac{\partial^\alpha\mu_{uuu}}{\partial t^\alpha}-v\frac{\partial^\alpha\mu_{uuv}}{\partial t^\alpha}+2g(x)\mu_{uv}+g(x)u\mu_{uuv}-\phi_{xuu}=0 \label{19uu}, \\
-u\frac{\partial^\alpha\mu_{uvv}}{\partial t^\alpha}-\frac{\partial^\alpha\mu_{vv}}{\partial t^\alpha}-v\frac{\partial^\alpha\mu_{vvv}}{\partial t^\alpha}+g(x)u\mu_{vvv}-\phi_{xvv}=0
\label{19vv} 
\end{align}
correspondingly. 

Next, by subtracting the product of $f(x)$ and \eqref{19vv} from \eqref{19uu}, and utilizing the identities \eqref{11u15v} and \eqref{f11v15u}, we obtain
\begin{equation}
2g(x)\mu_{uv} = 0.
\end{equation}
Given that $g(x) \neq 0$, it follows that $\mu_{uv} = 0$. Applying this result, equation \eqref{11v} reduces to $\phi_{vv} = 0$ and \eqref{15u} reduces to $\phi_{uu} = 0$. Consequently, $\phi(x, t, u, v)$ is established as a linear function with respect to the dependent variables $u$ and $v$.

In an analogous manner, we consider the second-order derivatives of \eqref{arav} with respect to the dependent variables. By subtracting the product of $f(x)$ and the second-order derivative of \eqref{arav} with respect to $v$ from its second-order derivative with respect to $u$, we obtain
\begin{equation}
2g(x)\phi_{uv} = 0.
\end{equation}
Since $g(x) \neq 0$, it follows that $\phi_{uv} = 0$. Utilizing the previously established relations in \eqref{11v} and \eqref{15u}, we find that $\mu_{vv} = 0$. This confirms that $\mu(x, t, u, v)$ is also a linear function with respect to $u$ and $v$. 


\end{proof}

\end{document}
